\documentclass[a4paper,oneside,12pt]{extarticle}

\usepackage[utf8]{inputenc}
\usepackage[english]{babel}
\usepackage{amssymb,amsthm,mathtools,enumitem,geometry,hyperref,booktabs}
\usepackage{graphicx}
\usepackage{xcolor}
\usepackage{tikz}
\usetikzlibrary{calc,patterns,decorations.pathreplacing}
\usepackage{tikz-cd}
\usepackage{pgfplots}
\usepgfplotslibrary{fillbetween}
\pgfplotsset{compat=1.18}
\usepackage{float}

\geometry{lmargin=15mm,rmargin=15mm,tmargin=10mm,bmargin=10mm,includeheadfoot}
\pagestyle{plain}
\frenchspacing

\theoremstyle{plain}
\newtheorem{proposition}{Proposition}
\newtheorem{lemma}[proposition]{Lemma}

\theoremstyle{definition}
\newtheorem{definition}[proposition]{Definition}
\newtheorem{example}[proposition]{Example}

\theoremstyle{remark}
\newtheorem{remark}[proposition]{Remark}

\newcommand*{\E}{\mathbb{E}}
\newcommand*{\Var}{\mathrm{Var}}
\newcommand*{\mP}{\mathbb{P}}
\newcommand*{\mR}{\mathbb{R}}
\newcommand*{\mbF}{\mathcal{F}}
\newcommand*{\ve}{\varepsilon}




\begin{document}
\begin{titlepage}
\begin{center}

{\footnotesize\scshape\color{gray}
Working Notes in Quantitative Risk and Modeling
}

\vfill

{\LARGE\bfseries Homoscedasticity and Its Testing:\\[4pt]
From Conditional Expectation to the Breusch--Pagan Test}

\vspace{36pt}

{\normalsize Boris Garbuzov}

\vfill

{\Large\bfseries Part~9~/~12}

\vspace{6pt}

{\Large\bfseries Sequential Sampling and the Bayesian Connection}

\vfill

{\small July 2026}

\end{center}
\end{titlepage}

\setcounter{tocdepth}{2}
\tableofcontents
\newpage
\section{Sequential Sampling of Random Quantities}
\label{sec:sequential_sampling}
%--------------------------------------------------------------------

\subsection{Motivation}

We routinely manipulate measure-dependent functions: we sample them,
integrate them, take moments, differentiate them, form derived
quantities.  The density $f_P$ of a measure $P$ is one such object.
A natural question: can we iterate this process --- plug a random
draw back into the density, treat the result as a new random
variable, and repeat?  This section explores three such
constructions and their limiting behaviour.

\subsection{Setup}

Fix a probability measure $P$ on $\mathbb{R}$ with Lebesgue density
$f = f_P \geq 0$, $\int f = 1$.  Draw $X_1 \sim P$.  Since
$f(x) \geq 0$ and $\int f = 1$, the function $f$ maps $\mathbb{R}$
into $[0, \sup f\,]$, so all iterates stay bounded above by
$\sup f < \infty$ (assuming the density is bounded, as for the
normal family).

\subsection{Three sequences}

\paragraph{Sequence $P_n$ (chained push-forward).}
Define $X_{n+1} = f_{P_n}(X_n)$, where $P_n$ is the distribution
of $X_n$ and $f_{P_n}$ is its density.  Each step is a
\emph{push-forward}: $P_{n+1} = (f_{P_n})_* P_n$.  The density
changes at every step together with the sampling distribution.

\paragraph{Sequence $Q_n$ (orbit of a point under $f$).}
Fix the original density $f = f_{P_1}$ throughout.  Set
\[
  X_{n+1} = f(X_n), \qquad X_1 \sim P.
\]
This is deterministic given $X_1$: the sequence $X_1, f(X_1),
f(f(X_1)), \ldots$ is the \emph{orbit} of the initial point $X_1$
under the iterated application of $f$.  The randomness enters only
through the initial draw; thereafter $Q_n$ is the push-forward of
$P$ under $f^{\circ n}$ (the $n$-fold composition of $f$ with itself).

\paragraph{Sequence $R_n$ (fixed sample point, varying density).}
Keep $X_1$ fixed and let the density vary.  Concretely, if $P_n$
has density $f_n$, form
\[
  R_n = f_n(X_1).
\]
Here $X_1$ is a fixed (random) argument and $f_n$ evolves.  This
is closer to evaluating a likelihood at a fixed data point under a
changing model.

\subsection{Limiting behaviour of sequence $Q_n$}

The sequence $Q_n$ is the most concrete: it reduces to the dynamics
of the map $f: \mathbb{R} \to [0,\sup f\,]$.

\paragraph{Contraction.}
After one step, every point is mapped into $[0,\sup f\,]$.  If $f$
is a contraction on this interval --- i.e.\ $|f'(x)| < 1$ for all
$x$ in the range --- then by the Banach fixed-point theorem the orbit
converges to the unique fixed point $x^*$ satisfying
\[
  f(x^*) = x^*.
\]

\paragraph{Fixed point for the normal density.}
For $f(x) = \frac{1}{\sqrt{2\pi}}e^{-x^2/2}$ the equation
$f(x^*) = x^*$ has a solution near $x^* \approx 0.3989$.
(At $x=0$: $f(0) = 1/\sqrt{2\pi} \approx 0.399 > 0$; at $x=1$:
$f(1) \approx 0.242 < 1$; by the intermediate value theorem a
crossing exists.)  Moreover $f'(x) = -x f(x)$, so
$|f'(x^*)| = x^* \cdot f(x^*) = (x^*)^2 \approx 0.159 < 1$:
the fixed point is attracting and the orbit converges to it for
almost every starting point $X_1$.

The limiting distribution of $Q_n$ therefore degenerates to a point
mass at $x^*$: starting from any $X_1 \sim P$, the sequence
$f(X_1), f(f(X_1)), \ldots$ converges almost surely to $x^*$.

\begin{remark}[Degenerate limit]
The limit measure is $\delta_{x^*}$, a Dirac mass.  It has no
Lebesgue density.  So even if every $Q_n$ has a density, the limit
can lose absolute continuity.
\end{remark}

\subsection{Limiting behaviour of sequence $P_n$}

The chained push-forward is more complex because the density
changes at each step.  Even in the normal case, $P_2$ is the
distribution of $f(X_1)$ where $X_1 \sim \mathcal{N}(0,1)$: its
density is obtained by the change-of-variables formula applied to
the non-monotone map $f$.  Since $f$ is even (symmetric about 0)
and has a unique maximum at 0, the map is two-to-one on
$\mathbb{R}\setminus\{0\}$ and the density of $P_2$ has support
in $(0, 1/\sqrt{2\pi}\,]$.

For the general sequence $P_n$, the densities $f_{P_n}$ contract
their support and, under mild conditions, the sequence converges
to a degenerate measure --- the rate and limit depending on $P$.

\subsection{Dependence on the original distribution $P$}

All three sequences depend only on $P$.  Different choices of $P$
produce qualitatively different behaviour:

\begin{itemize}
  \item \textbf{Light-tailed unimodal $P$} (e.g.\ normal): the
    density is bounded, the fixed point $x^*$ exists and is
    attracting, the $Q_n$ sequence converges to $\delta_{x^*}$.

  \item \textbf{Heavy-tailed $P$} (e.g.\ Cauchy): the density is
    still bounded ($\sup f = 1/\pi$), but the orbit dynamics depend
    on the fixed-point equation for the Cauchy density.

  \item \textbf{Multimodal $P$}: the density map $f$ may have
    multiple fixed points or even cycles, leading to more complex
    orbit structure (period-doubling, chaos).

  \item \textbf{Unbounded density} (e.g.\ Beta$(1/2, 1/2)$,
    which diverges at 0 and 1): $\sup f = \infty$ and the
    contraction argument breaks down; the sequence may diverge.
\end{itemize}

\subsection{Connection to Bayesian iteration and Markov chains}

The construction in sequence $P_n$ is reminiscent of Bayesian
updating: starting from a prior, each observation shifts the
measure.  But here the ``data'' is the density itself evaluated at
the current draw --- a self-referential structure.

More formally, $Q_n$ defines a (deterministic given $X_1$)
dynamical system on $\mathbb{R}$.  The randomness in $X_1$ induces
a distribution over orbits.  If $f$ has a unique attracting fixed
point $x^*$, the empirical measure of the orbit concentrates on
$\delta_{x^*}$ regardless of the starting distribution $P$
(as long as $P$ assigns positive mass to the basin of attraction
of $x^*$, which for the normal density is all of $\mathbb{R}$).

\begin{remark}[Open questions]
The limiting behaviour of $P_n$ (chained push-forward) for general
$P$ is not fully characterised here.  The key quantities are:
(a) whether $f_{P_n}$ has a fixed point; (b) the spectral radius
of the linearisation $Df_{P_n}$ at that fixed point; (c) whether
the push-forward operation preserves absolute continuity.
These depend on the smoothness and tail behaviour of $P$.
\end{remark}





\end{document}
